\documentclass[11pt]{article}
\usepackage[a4paper,margin=25mm]{geometry}
\usepackage{amsmath,amssymb,amsthm,hyperref}
\newtheorem{theorem}{Theorem}
\newtheorem{lemma}{Lemma}
\newcommand{\Q}{\mathbb Q}
\newcommand{\G}{\Gamma(1/4)}
\newcommand{\D}{\mathcal D}
\title{Coefficient cancellations in mixed Berndt integrals\\
\large MathDB 367221 --- Task 1}
\author{Principia Math in Progress}
\date{12 September 2026}
\begin{document}
\maketitle

Put $G=\G$ and define
\[
 I_\delta^\epsilon(s)=\int_0^\infty
 \frac{x^s(\sinh x+\epsilon\sin x)}
 {(\sinh^2x+\sin^2x)(\cosh x+\delta\cos x)}\,dx,
 \qquad \epsilon,\delta\in\{1,-1\}.
\]
The four coefficient pairs in Zhou's Theorems 5.1 and 5.4
\cite{Zhou} are specified by
\begin{align}
 I_+^-(4m)&=c_{1,m}\frac{G^{8m}}{\pi^{2m}}
                  +c_{2,m}\frac{G^{8m+4}}{\pi^{2m+2}},\label{c}\\
 I_-^+(4m)&=e_{1,m}\frac{G^{8m}}{\pi^{2m}}
                  +e_{2,m}\frac{G^{8m+4}}{\pi^{2m+2}},\label{e}\\
 I_+^+(4m-2)&=d_{1,m}\frac{G^{8m-4}}{\pi^{2m-1}}
                  +d_{2,m}\frac{G^{8m}}{\pi^{2m+1}},\label{d}\\
 I_-^-(4m-2)&=f_{1,m}\frac{G^{8m-4}}{\pi^{2m-1}}
                  +f_{2,m}\frac{G^{8m}}{\pi^{2m+1}}.\label{f}
\end{align}
All coefficients are rational. Equations \eqref{c}--\eqref{d} hold for
$m\geq1$; \eqref{f} holds for $m\geq2$. These are the source's
proved structural evaluations, not additional conjectural assumptions.

\begin{theorem}\label{main}
For every $m\geq1$, $c_{2,m}+e_{2,m}=0$. For every $m\geq2$,
$d_{1,m}+f_{1,m}=0$. At $m=1$, the second cancellation also holds
with $f_{1,1}$ interpreted as the coefficient of $G^4/\pi$
in the exceptional formula of Zhou's Remark 5.5.
\end{theorem}
This resolves Conjecture 5.8, with its first-index convention made
explicit. In particular, the exceptional formula is treated directly;
it is not an instance of \eqref{f}.

\section*{An alternating series and its symmetry}
For even $k\geq2$ and $y>0$, put
\[
 F_k(y)=\sum_{n=1}^\infty
 \frac{(-1)^{n-1}n^{k-1}}{\sinh(ny)}.
\]
This series and its first derivative converge locally uniformly for
$y>0$, since their summands are bounded by a polynomial in $n$
times $e^{-ny}$. For even $k\geq4$ write
\[
 E_k(\tau)=1-\frac{2k}{B_k}\sum_{n=1}^\infty
 \frac{n^{k-1}q^n}{1-q^n},\qquad q=e^{2\pi i\tau},
\]
where $B_k$ is the Bernoulli number. Define
\[
 H_k(\tau)=-E_k(\tau/2)+(1+2^k)E_k(\tau)-2^kE_k(2\tau).
\]
Splitting the Lambert series into even and odd indices gives
\begin{equation}
 F_k(y)=\frac{B_k}{k}H_k(iy/\pi).                 \label{lambert}
\end{equation}
For clarity, if $L(t)=\sum_{n\geq1}n^{k-1}t^n/(1-t^n)$,
the corresponding identity is
$F_k(y)=2\{L(e^{-y})-(1+2^k)L(e^{-2y})+2^kL(e^{-4y})\}$.

The usual modular transformation $E_k(-1/\tau)=\tau^kE_k(\tau)$
implies
\begin{equation}
 H_k(-1/\tau)=\tau^kH_k(\tau),\qquad
 F_k(\pi^2/y)=(-1)^{k/2}(y/\pi)^kF_k(y).          \label{symmetry}
\end{equation}
Indeed the terms at $\tau/2$ and $2\tau$ interchange, with factors
$2^k$ and $2^{-k}$. At $y=\pi$, \eqref{symmetry} gives
\begin{align}
 F_k'(\pi)&=-\frac{k}{2\pi}F_k(\pi)
                       &&(k\equiv0\pmod4),       \label{even}\\
 F_k(\pi)&=0            &&(k\equiv2\pmod4).       \label{odd}
\end{align}
The derivative identity follows by differentiating
\eqref{symmetry} before substituting $y=\pi$.

\section*{The two possible weights}
Set
\[
 A=\frac{G^4}{4\pi^3},\qquad
 \D=\frac{1}{2\pi i}\frac{d}{d\tau}.
\]
We use standard Eisenstein-series facts: $E_k$ is a weighted
homogeneous polynomial over $\Q$ in $E_4,E_6$, and
\begin{equation}
 \D E_k-\frac{k}{12}E_2E_k
 \text{ is a modular form of weight }k+2
 \text{ with rational Fourier coefficients}.       \label{serre}
\end{equation}
Thus the form in \eqref{serre} is also in $\Q[E_4,E_6]$.
The relevant special values are
\begin{equation}\label{cm}
\begin{array}{c|ccc}
 \tau&E_2(\tau)&E_4(\tau)&E_6(\tau)\\ \hline
 i/2&6/\pi-3A/2&33A^2/4&-189A^3/8\\
 i&3/\pi&3A^2/4&0\\
 2i&3/(2\pi)+3A/8&33A^2/64&189A^3/512 .
\end{array}
\end{equation}
These facts and values are recorded, with the derivative method,
in Wang \cite[Sections 2 and 4]{Wang}. They also follow from the
classical theta-constant formulas and the quadratic transformation.

Weighted homogeneity first gives
\begin{equation}
 H_k(i)\in\Q A^{k/2}.                                \label{value}
\end{equation}
By the chain rule,
\[
 \D H_k(i)=-\tfrac12(\D E_k)(i/2)
          +(1+2^k)(\D E_k)(i)-2^{k+1}(\D E_k)(2i).
\]
Inserting \eqref{serre} and \eqref{cm}, every term not involving $1/\pi$
belongs to $\Q A^{k/2+1}$. The $1/\pi$ contribution is precisely
$kH_k(i)/(4\pi)$: the three chain-rule factors multiply
$6/\pi,3/\pi,3/(2\pi)$ to the common value $3/\pi$.
Consequently
\begin{equation}
 \D H_k(i)-\frac{k}{4\pi}H_k(i)\in\Q A^{k/2+1}.
                                                               \label{derivative}
\end{equation}
Since $dH_k(iy/\pi)/dy=-2\D H_k(iy/\pi)$,
\eqref{lambert}--\eqref{derivative} prove
\begin{equation}
 \begin{cases}
 F_{4m}'(\pi)\in\Q A^{2m}/\pi,&m\geq1,\\
 F_{4m-2}'(\pi)\in\Q A^{2m},&m\geq2.
 \end{cases}                                                   \label{purity}
\end{equation}
For the first line use \eqref{even} and \eqref{value};
for the second use \eqref{odd} and \eqref{derivative}.

\section*{Return to the four integrals}
Let
\[
 \overline C_{s,2}(y)=\sum_{n\geq1}\frac{(-1)^{n-1}n^s}{\cosh^2(ny)},
 \qquad
 \overline S_{s,2}(y)=\sum_{n\geq1}\frac{(-1)^{n-1}n^s}{\sinh^2(ny)}.
\]
The elementary identity
\[
 \frac1{\cosh^2t}+\frac1{\sinh^2t}
       =\frac{4\cosh(2t)}{\sinh^2(2t)}
\]
and termwise differentiation give
\begin{equation}
 \overline C_{k,2}(\pi/2)+\overline S_{k,2}(\pi/2)
                  =-4F_k'(\pi).                     \label{sum}
\end{equation}
Zhou's contour evaluations
\cite[(5.15), (5.16), (5.33), (5.34)]{Zhou} yield
\begin{align}
 I_+^-(4m)+I_-^+(4m)
   &=\frac{(-1)^{m-1}\pi^{4m+1}}{2\,4^m}
       \bigl(\overline C_{4m,2}(\pi/2)+
                        \overline S_{4m,2}(\pi/2)\bigr),
                                                       \label{firstsum}\\
 I_+^+(4m-2)+I_-^-(4m-2)
   &=\frac{(-1)^{m-1}\pi^{4m-1}}{4^m}
       \bigl(\overline C_{4m-2,2}(\pi/2)+
                        \overline S_{4m-2,2}(\pi/2)\bigr)
                       \quad(m\geq2).                 \label{secondsum}
\end{align}
It follows from \eqref{purity} and \eqref{sum} that
\begin{align}
 I_+^-(4m)+I_-^+(4m)&\in\Q G^{8m}/\pi^{2m},\label{pure1}\\
 I_+^+(4m-2)+I_-^-(4m-2)&\in\Q G^{8m}/\pi^{2m+1}
                                                  \quad(m\geq2).\label{pure2}
\end{align}
The powers simplify as follows:
\[
 \frac{\pi^{4m+1}A^{2m}}{\pi}
       =\frac{G^{8m}}{4^{2m}\pi^{2m}},
 \qquad
 \pi^{4m-1}A^{2m}=\frac{G^{8m}}{4^{2m}\pi^{2m+1}}.
\]

The known algebraic independence of $G$ and $\pi$ implies the
linear independence over $\Q$ of each pair of monomials in
\eqref{c}--\eqref{f}; this is also explicitly recalled in
\cite[Section 3, after equations (101)--(103)]{XuZhao}.
Comparing \eqref{pure1} with \eqref{c}+\eqref{e} gives
$c_{2,m}+e_{2,m}=0$. Comparing \eqref{pure2} with
\eqref{d}+\eqref{f} gives $d_{1,m}+f_{1,m}=0$ for $m\geq2$.

Finally, Zhou's Example 5.3 and Remark 5.5 state
\[
 I_+^+(2)=\frac{G^8}{512\pi^3}-\frac{G^4}{64\pi},
 \qquad
 I_-^-(2)=\frac{G^4}{64\pi}-\frac{\pi}{8}.
\]
Thus $d_{1,1}=-1/64$ and $f_{1,1}=1/64$ in the stated convention.
The extra term $-\pi/8$ explains why \eqref{secondsum} was only
used for $m\geq2$. This completes the proof of Theorem \ref{main}.

\section*{Tasks 2--5}
The following tasks are directions for further work, not claims
proved in Task 1.
\begin{enumerate}
\setcounter{enumi}{1}
\item Derive recurrences or closed formulas for the surviving
coefficients $c_{1,m}+e_{1,m}$ and $d_{2,m}+f_{2,m}$ from the
polynomial expressions for Eisenstein series.
\item Determine the coefficient cancellations at rectangular
complex multiplication points, replacing the self-dual value
$y=\pi$ by $y=\pi\sqrt r$ with positive rational $r$.
\item Classify rational linear combinations of
$E_k(\tau/d)$ and $E_k(d\tau)$ whose derivative at an elliptic
fixed point has only one of the two possible weights.
\item Extend the cancellation argument to higher powers of the
hyperbolic denominators, keeping explicit account of the
exceptional small-index terms.
\end{enumerate}

\begin{thebibliography}{9}
\bibitem{Zhou} J. Zhou, \emph{Mixed Berndt-Type Integrals and
Generalized Barnes Multiple Zeta Functions},
\href{https://arxiv.org/abs/2502.15173v4}{arXiv:2502.15173v4},
3 February 2026. Journal DOI:
\href{https://doi.org/10.1016/j.aam.2025.103027}
{10.1016/j.aam.2025.103027}.
The proof uses the current author manuscript.
\bibitem{Wang} W. Wang, \emph{Evaluation of reciprocal sums of
hyperbolic functions using quasimodular forms},
\href{https://arxiv.org/abs/2312.16929v2}{arXiv:2312.16929v2},
31 March 2025; J. Math. Anal. Appl. 549 (2025), 129526.
\bibitem{XuZhao} C. Xu and J. Zhao, \emph{Reciprocal Hyperbolic
Series of Ramanujan Type}, Mathematics 12 (2024), 2974,
\href{https://doi.org/10.3390/math12192974}{doi:10.3390/math12192974}.
\end{thebibliography}
\end{document}
